% ============================================================================
%  Section VI — Results  (IOP MST)
%  Passive voice / third person throughout. Reporting only; interpretation and
%  limitations are in Section VII (discussion.tex).
%  Requires: booktabs, graphicx, pifont.
%  Figures (results/): fig_optimism_gap.png, fig_confusion_group.png,
%    fig_tradeoffs_main.png, fig_recipe_stacking.png.
\providecommand{\cmark}{\ding{51}}
\providecommand{\xmark}{\ding{55}}
% ============================================================================

\section{Results}
\label{sec:results}

\subsection{The Optimism Gap of Random-Split Evaluation}
\label{subsec:gap}
Table~\ref{tab:optimism} reports the classification accuracy of the three
baseline models under the four evaluation protocols, and Fig.~\ref{fig:gap}
visualises the same result. Under the RANDOM protocol every model attains
essentially perfect accuracy, reproducing---and slightly exceeding---the $98\%$
reported by the source data article's random-split baseline. Under the
leakage-free GROUP protocol the accuracy of the same features and models falls to
roughly one-half, an \emph{optimism gap} of approximately $41$ to $47$ accuracy
points. The CROSS-SPEED and CROSS-LOAD protocols, which additionally impose an
operating-condition shift, yield comparably reduced accuracy.

\subsubsection{Which mechanism produces the gap}
\label{subsec:decomp}
A gap of that size is usually attributed to two causes at once: overlapping
windows sharing samples across the partition boundary, and windows inheriting the
identity of the recording they were cut from. These are separable, and a third
candidate---simple temporal proximity between neighbouring windows---sits between
them. Table~\ref{tab:decomp} removes them one at a time. Discarding the $50\%$
overlap halves the window count and guarantees that no sample appears in both
partitions; holding out a contiguous temporal block within every recording
additionally removes near-neighbour proximity while still training on every
recording; and the recording-wise protocol finally removes recording identity as
well.

The attribution is unambiguous. Removing window overlap changes nothing: a random
split of strictly non-overlapping windows still classifies at $100.0\%$. Removing
temporal proximity as well changes nothing either, with the blocked protocol also
at $100.0\%$. The entire gap---$46.7$ points---appears at the last step, when
whole recordings are held out. Sharing samples between partitions is therefore not
the operative mechanism, contrary to the explanation usually offered. Every window
cut from one recording carries that recording's particular mounting, operating
point and session conditions, and this signature is apparently sufficient to
identify the recording, and hence its label, from any window of it. Halving the
windows or separating them in time leaves that signature intact. Only holding out
the recording removes it.

This has a direct methodological consequence. Guarding against leakage by
enforcing non-overlapping windows, a common precaution, would have offered no
protection whatever on this dataset, and the same is likely wherever each
operating condition is captured in a single continuous acquisition. The unit that
must be held out is the recording.

Two further observations follow from Table~\ref{tab:optimism}: the leakage-free
accuracy remains well above the $16.7\%$ chance level, and the lower-capacity
models (KNN, SVM) attain higher GROUP accuracy than the Random Forest.

The seed-averaged confusion matrix of the GROUP protocol
(Fig.~\ref{fig:confusion}) localises the residual difficulty. The healthy state
is the most reliably identified, followed by the supply-imbalance and
broken-rotor-bar faults, whereas the two lubrication-severity states (B1 and B2)
are frequently confused with each other and with the remaining bearing
conditions, in agreement with the $50$~Hz band limit established in
Section~\ref{sec:charact}. Because this per-class behaviour is markedly uneven,
accuracy alone overstates how well the fine bearing states are resolved: the
macro-averaged $F_1$-score of the same configuration is $51.1\%$, and per-class
recall falls from $0.87$ for the healthy state and $0.66$ for the supply
imbalance to $0.35$ for B3 and $0.16$ for B2. The lubrication-severity states
therefore dominate the residual error, and the headline figures below should be
read with that qualification.

\subsection{Accuracy--Resource Trade-offs}
\label{subsec:tradeoff}
Each sensing and communication parameter was varied in turn under the GROUP
protocol (Fig.~\ref{fig:tradeoffs}). Reducing the sampling rate did not degrade
accuracy: as shown in Table~\ref{tab:sampling}, every rate at or below $25$~Hz
matched or exceeded the accuracy obtained at $100$~Hz, with a maximum near
$25$~Hz. Reducing the channel set was likewise not detrimental---the three raw
axes alone outperformed the full six-channel set---while lengthening the window
from $2$~s to $4$~s increased accuracy at the cost of detection latency.

The combination of these frugal choices was then examined
(Table~\ref{tab:recipe} and Fig.~\ref{fig:recipe}). The resource-optimal
configuration---$25$~Hz sampling, the three raw channels, and a $4$~s
window---raised the Random Forest accuracy from $52.8\%$ at the naive
$100$~Hz/six-channel/$2$~s baseline to $75.9\%$, while reducing the raw data
volume by approximately eight-fold. Robustness to further deployment stresses was
also measured: reducing the amplitude resolution to $4$~bits lowered accuracy by
only a few points, injected packet loss of up to $20\%$ left accuracy essentially
unchanged, and the six-feature subset of the source baseline performed within
noise of the full feature set. The largest sensitivity was to sensor mounting: a
simulated $30^{\circ}$ change in orientation degraded accuracy substantially.

Table~\ref{tab:measured} re-estimates the endpoints of that recipe over ten
partitions rather than five and scores them per held-out recording. The small
shifts relative to Table~\ref{tab:recipe}---$52.8$ against $53.9$ for the
baseline, $75.9$ against $75.0$ for the frugal configuration---follow from the
larger number of partitions and are an order of magnitude smaller than the
dispersion reported below.

The dispersion of these figures deserves care. Estimates were repeated over ten
random recording-level partitions and scored separately on each held-out
recording, so that the reported interval reflects recording-to-recording
variation rather than the effect of re-partitioning the same data
(Section~\ref{subsec:stats}). Table~\ref{tab:measured} gives the result. The
intervals are wide---roughly $\pm12$ points---because performance varies
substantially from one recording to the next, and an interval computed over seeds
instead would have been about three times narrower and correspondingly
misleading.

Compared on their per-recording differences, the resource-frugal configuration
classifies the health state $21.1$ percentage points more accurately than the
baseline, and a two-sided paired permutation test over the $36$ recordings gives
$p=0.020$. The intermediate step tells a different story: reducing the sampling
rate from $100$ to $25$~Hz alone accounts for $9.4$ points, but that difference is
not significant at the recording level ($p=0.26$). The rate reduction on its own
is therefore not established as an improvement by these data, even though the
combination of rate reduction, channel selection and the longer window is. This
distinction is invisible to a test conducted over repeated re-partitionings, which
would have reported both differences as highly significant.

\subsection{What the sampling-rate effect actually is}
\label{subsec:antialias}
The sweep of Table~\ref{tab:sampling} shows accuracy rising as the sampling rate
falls, which invites the interpretation that the diagnostic information is
concentrated at low frequency and that discarding the rest helps. That
interpretation cannot be correct as stated. Shaft rotational frequencies in this
dataset lie between $13.1$ and $24.9$~Hz (Table~\ref{tab:speeds}), so decimation
to $25$~Hz, whose Nyquist frequency is $12.5$~Hz, removes every rotational
fundamental in the dataset---including the peaks identified in
Section~\ref{sec:charact}---and accuracy nevertheless improves.

Three effects are confounded in a naive rate sweep: the band that survives, the
number of samples the features are computed from, and the aliasing that occurs if
the anti-alias filter is omitted. Table~\ref{tab:antialias} separates them by
holding each constant in turn. Band-limiting alone is obtained by low-pass
filtering at $12.5$~Hz while retaining the $100$~Hz rate, so that the frequency
content matches the decimated condition but the window still contains $200$
samples; proper decimation to $25$~Hz then reduces the window to $50$ samples at
the same bandwidth; and repeating the decimation with the anti-alias filter
disabled isolates aliasing.

The outcome contradicts the simple reading. Restricting the band to $12.5$~Hz at
an unchanged sampling rate \emph{reduces} accuracy by $2.7$ points, so the content
above $12.5$~Hz is not merely useless---the full-rate configuration is using it.
The entire improvement, $12.8$ points, appears at the second step, where the only
change is that the same band-limited signal is represented by a quarter as many
samples per window. Disabling the anti-alias filter costs $24.9$ points, which
both confirms that the decimation is performed correctly and rules out any
suggestion that aliasing folds useful content into the observable band.

The mechanism is therefore an interaction between the sampling rate and the
feature set, not a property of the vibration signal alone. The $14$ time-domain
descriptors are not invariant to the sampling rate: zero-crossing rate, amplitude
entropy over a fixed $32$-bin histogram, and the higher-order moments are all
computed over $W$ samples, and at $100$~Hz a signal band-limited to $12.5$~Hz is
four times oversampled, so consecutive samples are strongly correlated and the
descriptors are estimated from far fewer effectively independent values than their
sample count suggests. Decimation removes that redundancy. The practical
consequence for a node is unchanged---$25$~Hz remains the better operating
point---but the reason is not that the fault information lies below $12.5$~Hz, and
any study that sweeps sampling rate with a fixed time-domain feature set measures
these two effects together rather than the effect of rate alone.

\subsection{Deployment Implications: Communication and Edge Footprint}
\label{subsec:edge}
Table~\ref{tab:linkbudget} compares the bitrate of three telemetry strategies
with the sustainable throughput of common low-power radios. Streaming the raw
six-channel signal requires $9.6$~kbps and transmitting the extracted feature
vector requires $2.69$~kbps; both are carried by Bluetooth Low Energy, Zigbee, or
NB-IoT, but neither fits within the duty-cycle-limited throughput of a LoRaWAN
uplink, which admits only the $40$~bps on-node decision stream.
Table~\ref{tab:footprint} lists the leakage-free accuracy and serialised size of
four candidate models: a $4.7$~kB logistic-regression classifier exceeded the
accuracy of a $1.09$~MB Random Forest, and the time-domain feature extraction
requires only $\mathcal{O}(WC)$ operations per window with no spectral transform.

% ---------------------------------------------------------------------------
\begin{table}[t]
\centering
\caption{Accuracy (\%) by evaluation protocol and model (five-seed means; $100$~Hz, six channels, $2$~s. Dispersion on the recording unit is reported in Tables~
ef{tab:decomp} and~
ef{tab:measured}).}
\label{tab:optimism}
\renewcommand{\arraystretch}{1.2}
\begin{tabular}{@{}lccc@{}}
\toprule
\textbf{Protocol} & \textbf{RF} & \textbf{KNN} & \textbf{SVM} \\
\midrule
RANDOM (leaky)            & $100.0$ & $100.0$ & $100.0$ \\
GROUP (recording-holdout) & $52.8$ & $59.0$ & $58.4$ \\
CROSS-SPEED               & $49.9$ & $48.2$ & $50.1$ \\
CROSS-LOAD                & $49.5$ & $40.3$ & $40.1$ \\
\bottomrule
\end{tabular}
\end{table}

% ---------------------------------------------------------------------------
% ---------------------------------------------------------------------------
\begin{table}[t]
\centering
\caption{Attribution of the optimism gap (Random Forest, five seeds; intervals are
$95\%$ clustered on recording, $n=36$). Each row removes one further source of
information shared between the training and test partitions.}
\label{tab:decomp}
\renewcommand{\arraystretch}{1.2}
\begin{tabular}{@{}lccc@{}}
\toprule
\textbf{Partitioning} & \textbf{Sources retained} & \textbf{Accuracy (\%)} & \textbf{Macro-$F_1$ (\%)} \\
\midrule
RANDOM, $50\%$ overlap    & samples, proximity, identity & $100.0$ & $100.0$ \\
RANDOM, no overlap        & proximity, identity          & $100.0$ & $100.0$ \\
BLOCKED within recording  & identity                     & $100.0$ & $100.0$ \\
GROUP, no overlap         & none                         & $53.3$ & $49.1$ \\
GROUP, $50\%$ overlap     & none                         & $52.8$ & $48.4$ \\
CROSS-SPEED               & none, $+$ speed shift        & $49.9$ & $41.5$ \\
CROSS-LOAD                & none, $+$ load shift         & $49.5$ & $44.5$ \\
\midrule
\multicolumn{4}{@{}l}{\textit{Attribution}} \\
\multicolumn{2}{@{}l}{shared samples (window overlap)} & $+0.0$ & \\
\multicolumn{2}{@{}l}{temporal proximity}             & $+0.0$ & \\
\multicolumn{2}{@{}l}{recording identity}             & $\mathbf{+46.7}$ & \\
\bottomrule
\end{tabular}
\end{table}

\begin{table}[t]
\centering
\caption{Sampling-rate sweep under the GROUP protocol (Random Forest, five-seed means). Lower rates match or exceed $100$~Hz.}
\label{tab:sampling}
\renewcommand{\arraystretch}{1.2}
\begin{tabular}{@{}ccc@{}}
\toprule
\textbf{Rate (Hz)} & \textbf{Raw bitrate (kbps)} & \textbf{Accuracy (\%)} \\
\midrule
100  & 9.6 & $52.8$ \\
50   & 4.8 & $39.2$ \\
25   & 2.4 & $\mathbf{62.9}$ \\
20   & 1.9 & $60.9$ \\
12.5 & 1.2 & $60.9$ \\
10   & 1.0 & $58.8$ \\
\bottomrule
\end{tabular}
\end{table}

% ---------------------------------------------------------------------------
\begin{table}[t]
\centering
\caption{Incremental frugal recipe under the GROUP protocol
(five-seed means). The resource-optimal node exceeds the naive
baseline.}
\label{tab:recipe}
\renewcommand{\arraystretch}{1.2}
\begin{tabular}{@{}lccc@{}}
\toprule
\textbf{Configuration} & \textbf{RF} & \textbf{KNN} & \textbf{LogReg} \\
\midrule
Baseline ($100$~Hz, $6$ ch, $2$~s) & $52.8$ & $59.0$ & $59.6$ \\
\;$+\,25$~Hz                        & $62.9$ & $46.5$ & $57.3$ \\
\;$+\,$raw $3$ ch                   & $65.6$ & $51.4$ & $52.5$ \\
\;$+\,4$~s window (optimal)         & $\mathbf{75.9}$ & $54.9$ & $60.3$ \\
\bottomrule
\end{tabular}
\end{table}

% ---------------------------------------------------------------------------
\begin{table}[t]
\centering
\caption{Leakage-free (GROUP) accuracy of the Random Forest, scored per held-out
recording over ten random recording-level partitions. Intervals are $95\%$
intervals for the mean, clustered on recording ($n=36$); comparisons are
two-sided paired permutation tests on the per-recording differences.}
\label{tab:measured}
\renewcommand{\arraystretch}{1.2}
\begin{tabular}{@{}lcc@{}}
\toprule
\textbf{Configuration} & \textbf{Accuracy (\%)} & \textbf{Macro-$F_1$ (\%)} \\
\midrule
Baseline ($100$~Hz, $6$ ch, $2$~s)   & $53.9 \pm 12.5$ & $49.1$ \\
$\quad25$~Hz, $6$ ch, $2$~s          & $63.2 \pm 11.6$ & $59.2$ \\
Frugal ($25$~Hz, $3$ raw ch, $4$~s)  & $\mathbf{75.0 \pm 11.4}$ & $\mathbf{71.1}$ \\
\midrule
\multicolumn{3}{@{}l}{\textit{Paired comparison on per-recording differences} ($n=36$)} \\
Frugal $-$ Baseline  & $+21.1$ & $p=0.020$ \\
$25$~Hz $-$ Baseline & $+9.4$  & $p=0.26$ (n.s.) \\
\bottomrule
\end{tabular}
\end{table}

% ---------------------------------------------------------------------------
% ---------------------------------------------------------------------------
\begin{table}[t]
\centering
\caption{Separating band-limiting, sample count and aliasing (GROUP protocol,
Random Forest, five seeds; intervals are $95\%$ clustered on recording, $n=36$).
Conditions (b) and (c) contain identical frequency content and differ only in the
number of samples per window.}
\label{tab:antialias}
\renewcommand{\arraystretch}{1.2}
\begin{tabular}{@{}llccc@{}}
\toprule
 & \textbf{Condition} & \textbf{Samples/win.} & \textbf{Accuracy (\%)} & \textbf{Macro-$F_1$ (\%)} \\
\midrule
(a) & $100$~Hz, unmodified           & $200$ & $52.8$ & $48.4$ \\
(b) & $100$~Hz, low-pass $12.5$~Hz   & $200$ & $50.1$ & $45.1$ \\
(c) & $25$~Hz, anti-alias decimation & $50$  & $\mathbf{62.9}$ & $\mathbf{58.6}$ \\
(d) & $25$~Hz, no anti-alias filter  & $50$  & $38.0$ & $32.2$ \\
\midrule
\multicolumn{5}{@{}l}{\textit{Attribution of the $+10.1$ point net change}} \\
\multicolumn{3}{@{}l}{band-limiting alone, (a)$\rightarrow$(b)}        & $-2.7$ & \\
\multicolumn{3}{@{}l}{sample count at fixed band, (b)$\rightarrow$(c)} & $\mathbf{+12.8}$ & \\
\multicolumn{3}{@{}l}{aliasing, (c)$\rightarrow$(d)}                   & $-24.9$ & \\
\bottomrule
\end{tabular}
\end{table}

\begin{table}[t]
\centering
\caption{Communication link budget: telemetry bitrate versus sustainable uplink
capacity. Bitrates assume a $1$~s reporting hop, \texttt{int16} raw samples and
\texttt{float32} features; the assumed sustainable application-level uplinks are
$300$~kbps (Bluetooth Low Energy $4.2$), $100$~kbps (Zigbee, IEEE~802.15.4),
$30$~kbps (NB-IoT) and $55$~bps (LoRaWAN, a spreading-factor-$7$ physical rate of
$5.47$~kbps throttled by a $1\%$ duty cycle). \cmark~= sustainable,
\xmark~= infeasible.}
\label{tab:linkbudget}
\renewcommand{\arraystretch}{1.2}
\setlength{\tabcolsep}{4pt}
\begin{tabular}{@{}lccccc@{}}
\toprule
\textbf{Telemetry stream} & \textbf{Bitrate} & \textbf{BLE} & \textbf{Zigbee} &
\textbf{NB-IoT} & \textbf{LoRaWAN} \\
\midrule
Raw ($6$ ch, int16)     & $9.6$~kbps  & \cmark & \cmark & \cmark & \xmark \\
Features ($84$/s)       & $2.69$~kbps & \cmark & \cmark & \cmark & \xmark \\
Features (min., $6$/s)  & $192$~bps   & \cmark & \cmark & \cmark & \xmark \\
On-node decision        & $40$~bps    & \cmark & \cmark & \cmark & \cmark \\
\bottomrule
\end{tabular}
\end{table}

% ---------------------------------------------------------------------------
\begin{table}[t]
\centering
\caption{Edge-model footprint versus leakage-free (GROUP) accuracy at the
baseline configuration ($100$~Hz, six channels, $2$~s; five-seed means). Serialised sizes are
architecture-determined.}
\label{tab:footprint}
\renewcommand{\arraystretch}{1.2}
\begin{tabular}{@{}lcc@{}}
\toprule
\textbf{Model} & \textbf{GROUP accuracy (\%)} & \textbf{Serialised size} \\
\midrule
Random Forest ($200$ trees)     & $52.8$ & $1.09$~MB \\
Random Forest ($50$, depth $8$) & $52.9$ & $214$~kB \\
Decision tree (depth $8$)       & $44.3$ & $4.3$~kB \\
Logistic regression             & $\mathbf{59.6}$ & $4.7$~kB \\
\bottomrule
\end{tabular}
\end{table}

% ---------------------------------------------------------------------------
\begin{figure}[t]
\centering
\includegraphics[width=\columnwidth]{fig_optimism_gap.png}
\caption{Accuracy by evaluation protocol and model. Near-perfect accuracy under
the RANDOM split collapses once evaluation is leakage-free (GROUP) or imposes an
operating-condition shift (CROSS-SPEED, CROSS-LOAD).}
\label{fig:gap}
\end{figure}

% ---------------------------------------------------------------------------
\begin{figure}[t]
\centering
\includegraphics[width=\columnwidth]{fig_confusion_group.png}
\caption{Seed-averaged confusion matrix of the Random Forest under the GROUP
protocol (row-normalised; baseline configuration, five recording-level
partitions). The healthy state (H) is the most separable, followed by the
supply-imbalance (V) and broken-rotor-bar (R) faults; the lubrication-severity
states (B1, B2) are the most confused.}
\label{fig:confusion}
\end{figure}

% ---------------------------------------------------------------------------
\begin{figure}[t]
\centering
\includegraphics[width=\columnwidth]{fig_tradeoffs_main.png}
\caption{Accuracy--resource trade-offs under the GROUP protocol: sampling rate
(left), window length (centre), and channel subset (right). The dashed line marks
the $16.7\%$ chance level.}
\label{fig:tradeoffs}
\end{figure}

% ---------------------------------------------------------------------------
\begin{figure}[t]
\centering
\includegraphics[width=\columnwidth]{fig_recipe_stacking.png}
\caption{Incremental frugal recipe under the GROUP protocol. Combining the
resource-frugal choices raises Random Forest accuracy from $52.8\%$ to $75.9\%$
while reducing the raw data volume roughly eight-fold.}
\label{fig:recipe}
\end{figure}
